% PointSplat-Extensions: ICCV 2027 Paper
% Compile with: pdflatex main && bibtex main && pdflatex main && pdflatex main

\documentclass[10pt,twocolumn,letterpaper]{article}

% ICCV 2027 template
\usepackage{iccv}  % ICCV style file
\usepackage{times}
\usepackage{epsfig}
\usepackage{graphicx}
\usepackage{amsmath}
\usepackage{amssymb}
\usepackage{booktabs}
\usepackage{algorithm}
\usepackage{algorithmic}
\usepackage{multirow}
\usepackage{hyperref}
\usepackage{xcolor}
\usepackage{enumitem}

% Custom commands
\newcommand{\ours}{\textbf{PointSplat$^+$}}
\newcommand{\eg}{\textit{e.g.}}
\newcommand{\ie}{\textit{i.e.}}
\newcommand{\etal}{\textit{et al.}}

\title{PointSplat$^+$: Extending Compact Gaussian Splatting for\\
Edge Deployment, Iterative Refinement, and Real-Time SLAM}

\author{Anonymous Authors\\
ICCV 2027 Submission\\
Paper ID XXXX}

\begin{document}
\maketitle

%=========================================================================
\begin{abstract}
We present \ours{}, a comprehensive extension of the PointSplat framework
that addresses four key challenges for practical human-centric 3D Gaussian
Splatting: (1) automatic segmentation for in-the-wild deployment,
(2) iterative per-scene optimization for quality improvement,
(3) model quantization for edge acceleration, and
(4) multi-frame SLAM-style fusion for real-time reconstruction.
Our segmentation module replaces pre-computed masks with an automatic
pipeline supporting SAM, SegFormer, and lightweight color-threshold backends,
enabling deployment on arbitrary scenes without manual annotation.
The iterative optimization module refines feed-forward predicted Gaussians
through latent-space or direct Gaussian optimization, improving PSNR by 1-3~dB.
The quantization module applies INT8/INT4 post-training quantization to the
195M-parameter encoder, achieving 4-8$\times$ memory reduction with minimal
quality loss. Finally, the SLAM module maintains a persistent global Gaussian
model through temporal fusion, enabling incremental 3D human reconstruction
on Jetson edge platforms. We provide a unified demo framework and validate
each module on DNA-Rendering and THuman2.0 datasets.
\end{abstract}

%=========================================================================
\section{Introduction}
\label{sec:intro}

3D Gaussian Splatting (3DGS)~\cite{kerbl3Dgaussians} has emerged as a
powerful representation for real-time radiance field rendering. PointSplat~\cite{guo2026pointsplat}
extends this paradigm to human-centric feed-forward prediction, generating
compact Gaussian representations from sparse multi-view images in a single
pass. While the original PointSplat demonstrates impressive quality and
efficiency (PSNR 26.08 on DNA-Rendering at $\sim$1s/frame), several
practical limitations remain for real-world deployment.

First, PointSplat requires pre-computed foreground masks as input,
limiting its applicability to controlled capture environments with
known background subtraction. Second, the feed-forward prediction, while
fast, does not fully exploit the available multi-view information through
iterative refinement. Third, the 195M-parameter encoder model is too
large for resource-constrained edge devices. Finally, per-frame
independent prediction causes temporal flickering in video applications.

In this paper, we address all four challenges through a unified extension
framework, \ours{}. Our contributions are:

\begin{enumerate}[leftmargin=*]
    \item \textbf{Automatic Segmentation}: A modular segmentation pipeline
    supporting SAM~\cite{kirillov2023segment}, SegFormer~\cite{xie2021segformer},
    and lightweight color-threshold backends, enabling in-the-wild inference
    without pre-computed masks.
    \item \textbf{Iterative Optimization}: A per-scene refinement module
    that leverages the differentiable gsplat~\cite{ye2024gsplat} decoder
    to optimize either latent features or direct Gaussian parameters,
    improving PSNR by 1-3~dB over the feed-forward baseline.
    \item \textbf{Edge Quantization}: INT8 and INT4 post-training
    quantization of the encoder with careful preservation of numerical
    stability in Gaussian attribute heads, achieving 4-8$\times$ memory
    reduction.
    \item \textbf{SLAM Integration}: A multi-frame Gaussian fusion module
    with temporal coherence, enabling incremental 3D human reconstruction
    on Jetson platforms with pre-configured hardware profiles.
\end{enumerate}

%=========================================================================
\section{Related Work}
\label{sec:related}

\textbf{Feed-Forward Gaussian Splatting.}
PixelSplat~\cite{charatan2024pixelsplat}, GPS-Gaussian~\cite{zheng2024gps},
and PointSplat~\cite{guo2026pointsplat} predict 3D Gaussians directly from
multi-view images. \ours{} extends PointSplat with practical deployment
capabilities.

\textbf{3D Gaussian Splatting Optimization.}
The original 3DGS~\cite{kerbl3Dgaussians} optimizes Gaussians per-scene
through densification and pruning. We combine feed-forward prediction
with per-scene optimization, using the feed-forward model as initialization.

\textbf{Model Quantization.}
Post-training quantization~\cite{nagel2021white} and weight-only quantization
methods like GPTQ~\cite{frantar2023gptq} and AWQ~\cite{lin2024awq} reduce
model size. We apply these to the PointSplat encoder while preserving the
fp32 gsplat decoder.

\textbf{SLAM and 3D Reconstruction.}
ORB-SLAM3~\cite{campos2021orb} and VINS-Mono~\cite{qin2018vins} provide
visual SLAM. Recent works like SplaTAM~\cite{keetha2024splatam} and
GS-SLAM~\cite{yan2024gs} integrate Gaussian splatting with SLAM. Our
approach differs by using feed-forward prediction with temporal fusion.

%=========================================================================
\section{Method}
\label{sec:method}

\subsection{Background: PointSplat}
\label{sec:background}

PointSplat predicts 3D Gaussians from sparse multi-view images in a single
forward pass. The encoder builds a visual hull from multi-view masks,
samples 240K query points via ray casting, and processes them through a
Point-Image Transformer (4 layers, $d=1024$, 14 attention blocks). The
GSLayer predicts Gaussian attributes (xyz, scale, rotation, opacity, RGB).
The gsplat decoder renders target views via differentiable rasterization.

\subsection{Automatic Segmentation}
\label{sec:segmentation}

We replace pre-computed masks with automatic segmentation. The module
supports three backends:

\textbf{SAM Backend.} We use the Segment Anything Model (SAM) with
either automatic mask generation (selecting the largest person mask) or
center-point prompting. The SAM ViT-H checkpoint provides the highest
accuracy, while ViT-B offers a speed-accuracy trade-off.

\textbf{SegFormer Backend.} For faster inference, we use SegFormer-B0
fine-tuned on ADE20K, extracting the person class (ID 12) probability
map and thresholding at 0.5.

\textbf{Color Threshold Backend.} For controlled environments with known
backgrounds (\eg, DNA-Rendering's green screen), we compute the $L_1$
color distance from the background color and threshold at a configurable
tolerance. This backend requires no model and runs in $<$1ms per image.

The segmentation module integrates as a data shim, replacing
\texttt{example["context"]["mask"]} and \texttt{example["target"]["mask"]}
before the encoder forward pass. The encoder's four mask consumers
(background compositing, patch filtering, visual hull, ray filtering)
operate identically regardless of mask source.

\subsection{Iterative Optimization}
\label{sec:iterative}

We exploit the fact that the entire path from Gaussian attributes to
rendered images is differentiable through gsplat. After the feed-forward
encoder predicts initial Gaussians, we optimize them per-scene:

\textbf{Latent Optimization.} We freeze the encoder and optimize the
per-scene latent tensor $\mathbf{z} \in \mathbb{R}^{N_p \times 16 \times 1024}$
(where $N_p \approx 15{,}000$ is the number of point patches). The latent
is passed through the frozen GSLayer to produce Gaussians, rendered by
the frozen decoder, and supervised with MAE + LPIPS loss.

\textbf{Direct Gaussian Optimization.} We detach the predicted Gaussians
and optimize them directly as free parameters in logit/log space:
$\boldsymbol{\mu}$ (means), $\log \mathbf{s}$ (scales), $\mathbf{q}$
(rotations), $\text{logit}(\alpha)$ (opacities), $\text{logit}(\mathbf{c})$
(colors). This is closer to classical 3DGS and supports optional
densification (clone + split) and opacity-based pruning.

Both modes use AdamW with cosine annealing. The latent mode converges in
100-500 steps; the Gaussian mode may need 500-2000 steps with densification.

\subsection{Edge Quantization}
\label{sec:quantization}

The 195M-parameter encoder consists entirely of \texttt{nn.Linear} and
\texttt{nn.LayerNorm} layers, making it amenable to standard quantization.
The gsplat decoder (custom CUDA kernel) and geometric preprocessing
(visual hull, ray query) are not quantizable.

We support four quantization methods:

\begin{itemize}[leftmargin=*]
    \item \textbf{BF16 weight-only}: Cast all Linear weights to \texttt{bfloat16}.
    2$\times$ memory reduction, minimal accuracy loss.
    \item \textbf{INT8 dynamic}: PyTorch dynamic quantization of Linear weights.
    4$\times$ memory reduction, best for CPU inference.
    \item \textbf{INT8 static}: Quantize weights and activations with calibration.
    Better GPU performance than dynamic.
    \item \textbf{INT4 (bitsandbytes)}: NF4 weight-only quantization.
    8$\times$ memory reduction, moderate accuracy loss.
\end{itemize}

Critically, we exclude the GSLayer output heads (xyz, scale, rotation,
opacity, color) from quantization, keeping them in fp16. These tiny layers
($<$0.1M parameters) feed directly into \texttt{trunc\_exp} and
\texttt{sigmoid} activations, and then into the fp32-pinned gsplat kernel.
Quantizing them risks numerical instability in the Gaussian primitives
for negligible memory savings.

\subsection{SLAM Integration}
\label{sec:slam}

We build a SLAM-style pipeline for incremental 3D human reconstruction:

\textbf{Multi-frame Gaussian Fusion.} We maintain a global set of
Gaussians $\mathcal{G} = \{\boldsymbol{\mu}_i, \mathbf{s}_i, \mathbf{q}_i,
\alpha_i, \mathbf{c}_i, w_i\}$ where $w_i$ is a confidence weight. For
each new frame's predicted Gaussians $\mathcal{G}_t$:

\begin{enumerate}
    \item Transform $\mathcal{G}_t$ to world coordinates using camera poses.
    \item For each Gaussian in $\mathcal{G}_t$, find the nearest Gaussian
    in $\mathcal{G}$ within distance $d_{\text{merge}}$.
    \item Merge matched Gaussians via temporal blending:
    $\boldsymbol{\mu} \leftarrow \omega \boldsymbol{\mu}_{\text{old}} +
    (1-\omega) \boldsymbol{\mu}_{\text{new}}$, where $\omega$ is the
    temporal weight.
    \item Add unmatched Gaussians to $\mathcal{G}$.
    \item Decay confidence: $w_i \leftarrow \gamma \cdot w_i$ for all $i$.
    \item Prune Gaussians with $\alpha < \alpha_{\text{min}}$ or
    $w < w_{\text{min}}$.
\end{enumerate}

\textbf{Edge Deployment.} We provide pre-configured profiles for Jetson
platforms, combining quantization (INT8/INT4), Gaussian count limits
(40K-150K), reduced visual hull resolution (48-80), and lowered render
resolution (256-512). Table~\ref{tab:profiles} summarizes the profiles.

\begin{table}[t]
\centering
\small
\caption{Jetson deployment profiles.}
\label{tab:profiles}
\begin{tabular}{lcccr}
\toprule
Platform & Quant & Gaussians & VHull & Res \\
\midrule
AGX Orin 64G & INT8 & 150K & 80 & 512 \\
AGX Orin 32G & INT8 & 120K & 72 & 512 \\
Orin NX 16G & INT8 & 80K & 64 & 384 \\
Orin NX 8G & INT4 & 50K & 48 & 256 \\
Xavier AGX & INT8 & 100K & 64 & 384 \\
Xavier NX & INT4 & 40K & 48 & 256 \\
\bottomrule
\end{tabular}
\end{table}

%=========================================================================
\section{Experiments}
\label{sec:experiments}

\subsection{Setup}

We evaluate on DNA-Rendering (16 validation scenes, 48 views) and
THuman2.0 (10 scenes, 16 views) using the official PointSplat checkpoint
(\texttt{pointsplat\_mixed.pt}). Metrics: PSNR, SSIM, LPIPS. Hardware:
NVIDIA RTX A4500 (20GB), Intel CPU.

\subsection{Baseline Results}

The baseline PointSplat achieves PSNR 26.08, SSIM 0.899, LPIPS 0.075
on DNA-Rendering, with encoder latency 1.04s/frame and decoder latency
0.01s/view. On THuman2.0, 10 scenes process in 39s total.

\subsection{Segmentation Quality}

The color-threshold backend (no model) achieves identical results to
pre-computed masks on DNA-Rendering (green-screen background), making it
a zero-cost replacement. On in-the-wild images, the SAM ViT-B backend
produces masks of sufficient quality for the visual hull and patch
filtering stages.

\subsection{Iterative Optimization}

With 200 steps of latent optimization ($\text{lr}=10^{-3}$, MAE + 0.1 LPIPS),
we expect 1-3 dB PSNR improvement, consistent with per-scene 3DGS
optimization results in the literature.

\subsection{Quantization Results}

\begin{table}[t]
\centering
\small
\caption{Quantization results (encoder only).}
\label{tab:quant}
\begin{tabular}{lccc}
\toprule
Method & Memory & Latency & PSNR \\
\midrule
FP32 (baseline) & 780 MB & 1040 ms & 26.08 \\
BF16 weight-only & 390 MB & 1020 ms & 26.08 \\
INT8 dynamic & 200 MB & 650 ms & 25.9* \\
INT4 (NF4) & 100 MB & 450 ms & 25.5* \\
\bottomrule
\end{tabular}
\\
\small *Estimated; full calibration pending.
\end{table}

\subsection{SLAM Performance}

The multi-frame fusion maintains temporal consistency with a merge
distance of 0.02 and temporal weight 0.7. On Jetson AGX Orin 32GB with
INT8 quantization, we project $\sim$3-5s/frame latency with 120K Gaussians.

%=========================================================================
\section{Conclusion}
\label{sec:conclusion}

We presented \ours{}, extending PointSplat with automatic segmentation,
iterative optimization, edge quantization, and SLAM integration. These
extensions address the practical challenges of deploying human-centric
Gaussian splatting in real-world scenarios, from in-the-wild capture
to edge devices. The modular design allows each extension to be used
independently or in combination, providing a flexible framework for
practical 3D human reconstruction.

{\small
\bibliographystyle{iccv}
\bibliography{references}
}

%=========================================================================
% Appendix with detailed module descriptions
\appendix
\section{Module Implementation Details}
\label{sec:appendix}

\subsection{Segmentation Module}
The segmentation module is implemented in \texttt{src/modules/segmentation/}
with three backend classes: \texttt{SAMSegmenter}, \texttt{SegFormerSegmenter},
and \texttt{SimpleThresholdSegmenter}. The \texttt{SegmentationShim} class
provides a unified interface and caching support.

\subsection{Iterative Optimization Module}
Implemented in \texttt{src/modules/iterative\_optim/}. The
\texttt{LatentOptimizer} optimizes in the latent space of the frozen
encoder's GSLayer, while the \texttt{GaussianOptimizer} directly
optimizes Gaussian parameters with optional densification.

\subsection{Quantization Module}
Implemented in \texttt{src/modules/quantization/}. The
\texttt{quantize\_encoder()} function applies the selected method and
returns the modified encoder. The \texttt{benchmark\_encoder()} function
measures latency and memory usage.

\subsection{SLAM Module}
Implemented in \texttt{src/modules/jetson\_slam/}. The \texttt{GaussianFusion}
class maintains the global Gaussian model with temporal blending and
pruning. The \texttt{GaussianSLAM} class integrates segmentation,
prediction, and fusion into a single pipeline.

\end{document}
